% !TeX root = ../../../poisson-cubic.tex
\section{Fermat型と相対モジュライの既約成分}
\label{sec:pg-fermat-relative-components}

本節でいう既約成分は基礎位相空間の既約成分である。
$N=n+2\ge5$、$V=\mathbb C^N$とし、
\begin{equation}
 F_N=\sum_{i=1}^Nz_i^3,\qquad
 X_N=V(F_N),\qquad
 \mathfrak P_N^{\mathrm{Fer}}=\operatorname{Pois}(X_N)
\end{equation}
と置く。前節の相対Poissonスキームを$\mathscr P_n$、その非零部分を
$\mathscr P_n^\times$と書く。
$\left\{\begin{smallmatrix}N\\k\end{smallmatrix}\right\}_{\ge3}$は、
$\{1,\ldots,N\}$を大きさ$3$以上の$k$個の部分に分ける非順序集合分割の数とする。

\begin{theorem}\label{thm:pg-fermat-relative-count}
次が成り立つ。
\begin{enumerate}
\item $\mathfrak P_N^{\mathrm{Fer}}$の既約成分は、
各部分の大きさが$3$以上の集合分割で添字付けられる。$k$個の部分を持つ成分は
次元$3N-8k$、一般Poisson階数$2k$であり、
\begin{equation}\label{eq:pg-fermat-count}
 \#\operatorname{Irr}(\mathfrak P_N^{\mathrm{Fer}})
 =\sum_{k=1}^{\lfloor N/3\rfloor}
 \left\{\begin{smallmatrix}N\\k\end{smallmatrix}\right\}_{\ge3}.
\end{equation}
\item 固定モジュライスタック$\mathfrak{Pois}(X_N)$の既約成分は、
$N$の各部分が$3$以上の整数分割で添字付けられる。従って
\begin{equation}\label{eq:pg-fermat-stack-count}
 \#\operatorname{Irr}(\mathfrak{Pois}(X_N))
 =p_{\ge3}(N)=[q^N]\prod_{b\ge3}(1-q^b)^{-1}.
\end{equation}
\item 相対モジュライスタックには一般Poisson階数$2k$ごとに一つの既約成分があり、
\begin{equation}\label{eq:pg-relative-count}
 \#\operatorname{Irr}(\mathfrak{Pois}^{\mathrm{rel}}_n)=\lfloor N/3\rfloor.
\end{equation}
その次元は、$m=N-3k$として
\begin{equation}\label{eq:pg-relative-dim}
 d_{n,k}=k(3N-8)+\binom{m+2}{3}
 +k\left(10+6m+3\binom{m+1}{2}\right)-N^2
\end{equation}
である。
\end{enumerate}
零Poisson構造を除いても、これらの成分数は変わらない。
\end{theorem}

\subsection{固定Fermat超曲面}
\begin{proof}[第一項の証明]
Fermat式のHessianは$6\operatorname{diag}(z_1,\ldots,z_N)$であるから、
補題\ref{lem:pg-linearize}よりPoisson条件は
\begin{equation}\label{eq:pg-fermat-contractions}
 (\iota_{dz_i}P)\wedge(\iota_{dz_i}P)=0
 \qquad(1\le i\le N)
\end{equation}
を含む。定理\ref{thm:pg-normal}による分解
$P=\sum_{s=1}^kP_s$、$0\ne P_s\in\bigwedge^3W_s$を取る。
異なる$s,t$の積
$(\iota_{dz_i}P_s)\wedge(\iota_{dz_i}P_t)$は
$\bigwedge^2W_s\otimes\bigwedge^2W_t$という相異なる直和成分に属する。
従って\eqref{eq:pg-fermat-contractions}から、
各$dz_i$が非零となる$W_s$は高々一つである。
集合
\begin{equation}
 B_s=\{i:dz_i|_{W_s}\ne0\},\qquad
 V_B=\langle e_i:i\in B\rangle
\end{equation}
は互いに素で、$|B_s|\ge3$、$W_s\subset V_{B_s}$を満たす。
未使用の添字をいずれかの$B_s$に加えれば、$\{1,\ldots,N\}$の分割を得る。

逆に、そのような分割$\mathcal B$に対して
\begin{equation}
 Z_{\mathcal B}=\sum_{B\in\mathcal B}
 \widehat{\operatorname{Gr}}(3,V_B)\subset\bigwedge^3V
\end{equation}
と置く。各錐は独立な線形部分空間$\bigwedge^3V_B$に入るので、
この和は錐の直積であり閉既約である。
各点$P=\sum_BP_B$について、非零$P_B$の支持を$V_B$の基底に延長する。
この座標では$F_N=\sum_B\sum_{i\in B}z_i^3$は異なるブロックを混ぜず、
正規形定理の逆向きにより$P$はPoissonである。
従って有限個の$Z_{\mathcal B}$が$\mathfrak P_N^{\mathrm{Fer}}$の被約化を覆う。
各錐の次元は$3(|B|-3)+1$なので
\begin{equation}
 \dim Z_{\mathcal B}=3N-8|\mathcal B|.
\end{equation}

これらの閉既約部分集合の間に真の包含がないことを示す。
縮約写像$\bigwedge^2V^*\to V$の一般階数は
$Z_{\mathcal B}$上で$3|\mathcal B|$である。
階数が一定値以下となる条件は閉じているので、部分数が少ないものは多いものを含まない。
逆向きの包含は上の次元式から不可能である。
部分数が等しい二つの分割を比較するには、各支持$W_s$が
$B_s$の全座標方向に非零射影を持つ一般点を取る。
補題\ref{ne:blocks}により$P$から$W_s$の組が復元され、
従ってその座標支持$B_s$も復元される。
同じ部分数の別の分割へ入る余地はない。
従って各$Z_{\mathcal B}$が一つの既約成分である。
一般点のPoisson階数は正規形定理より$2|\mathcal B|$となる。
全ての錐は原点を含むので、原点は追加の成分を与えない。
\end{proof}

\begin{proof}[第二項の証明]
命題\ref{prop:pg-fixed-stack-algebraic}により成分の自己同型群軌道を数えればよい。
自己同型は$F_N(Az)=cF_N(z)$を満たす射影線形変換である。
Hessian行列式と連鎖律から
\begin{equation}\label{eq:pg-fermat-hessian-discriminant}
 \det\operatorname{Hess}(F_N)=6^N\prod_i z_i,
\end{equation}
\begin{equation}\label{eq:pg-fermat-hessian-transform}
 \prod_i(Az)_i=c^N\det(A)^{-2}\prod_i z_i
\end{equation}
を得る。一意分解性から$A$はmonomial行列であり、
$F_N(Az)=cF_N(z)$へ戻せば、スカラー行列を除いた対角成分は3乗根である。
従って
\begin{equation}\label{eq:pg-fermat-automorphism-group}
 \operatorname{Aut}(X_N)\simeq(\mu_3^N/\mu_3)\rtimes\mathfrak S_N.
\end{equation}
対角部分は各座標支持を保ち、$\mathfrak S_N$はその添字を置換する。
よって二つの集合分割が同じ軌道に属する条件は、部分の大きさの多重集合が等しいことである。
これは各部分が$3$以上の整数分割にほかならず、第二項を得る。
\end{proof}

\subsection{相対モジュライ}
\begin{proof}[第三項の証明]
一般Poisson階数$2k$の被約層を$\mathscr P_{n,k}$と書く。
正規形定理により、これは縮約写像$\bigwedge^2V^*\to V$の階数が$3k$となる
局所閉条件で定まる。まず3次式の非零代表を選んだ空間に移り、
\begin{equation}
 P_k=e_{123}+e_{456}+\cdots+e_{3k-2,3k-1,3k},\qquad
 \mathcal O_{N,k}=\mathrm{GL}(V)\cdot P_k
\end{equation}
と置く。$P_k$の支持を$W=\bigoplus_sW_s$、$V=W\oplus U$、$\dim U=m$とする。
補題\ref{ne:blocks}により、その安定化子$H_k$は$W_s$を置換する。
単位成分の元は各$W_s$上で$\mathrm{SL}(W_s)$に属し、
$U$上の可逆変換と$U\to W$の写像は任意である。
従って
\begin{equation}
 \dim H_k=8k+m^2+3km,\qquad
 \dim\mathcal O_{N,k}=k(3N-8).
\end{equation}

$P=P_k$に対するHessian条件\eqref{eq:pg-K}は、
異なる$W_s,W_t$間の混合2階微分の消滅と同値である。
実際、二つずつの共変ベクトルを各ブロックから選ぶとその混合成分が得られ、
逆向きはブロックごとのJacobian括弧から従う。
従って許される3次式の線形空間$\mathcal F_{N,k}$は
\begin{equation}
 F=c(u)+\sum_{s=1}^k
 \bigl(f_{s,3}(x_s)+f_{s,2}(x_s,u)+f_{s,1}(x_s,u)\bigr)
\end{equation}
で与えられる。ここで$f_{s,j}$の$x_s$次数は$j$、$u$次数は$3-j$であり、
\begin{equation}
 \ell_{N,k}:=\dim\mathcal F_{N,k}
 =\binom{m+2}{3}+k\left(10+6m+3\binom{m+1}{2}\right).
\end{equation}
滑らかな3次式の部分$\mathcal F_{N,k}^{\mathrm{sm}}$はFermat式を含む非空開集合である。
軌道上でHessian条件の線形写像は一定階数を持つので、その核は
$\mathcal O_{N,k}$上の階数$\ell_{N,k}$のベクトル束となる。
ここで$g\cdot F=F\circ g^{-1}$、$g\cdot P=(\bigwedge^3g)P$とする。
随伴束では$(gh,F_0)\sim(g,h\cdot F_0)$、$h\in H_k$と同一視する。
3次代表を選んだ被約階数層との比較射は
\begin{equation}
 \mathrm{GL}(V)\times^{H_k}\mathcal F_{N,k}^{\mathrm{sm}}
 \longrightarrow\widetilde{\mathscr P}_{N,k},\qquad
 [g,F_0]\longmapsto(g\cdot F_0,g\cdot P_k)
\end{equation}
である。$\mathrm{GL}(V)\to\mathcal O_{N,k}$のエタール局所での持上げを選べば逆射が得られ、
二つの持上げの差は$H_k$の作用なので逆射は降下する。
従ってこれは上の核ベクトル束の開部分との同型であり、
階数層は既約で次元$k(3N-8)+\ell_{N,k}$を持つ。

3次代表の変更は自由な作用
\begin{equation}
 a\cdot(F,P)=(aF,a^{-1}P),\qquad a\in\mathbb G_m,
\end{equation}
であり、$P\otimes F\in\bigwedge^3V\otimes L$を不変に保つ。
従って$\mathscr P_{n,k}$は既約で
$\dim\mathscr P_{n,k}=k(3N-8)+\ell_{N,k}-1$である。
$\mathrm{PGL}(V)$による商で$N^2-1$を引くと
\eqref{eq:pg-relative-dim}を得る。

各階数層の閉包が成分であることを示す。
縮約写像の階数条件から、小さい$k$の層の閉包は大きい$k$の層を含まない。
一方、隣り合う層の次元差は
\begin{equation}
 \dim\mathscr P_{n,k}-\dim\mathscr P_{n,k+1}
 =3(3k-1)(N-3k)+8>0
\end{equation}
なので、逆向きの包含もない。
零切断は、各$F$上の分解可能3重ベクトルをスカラー倍することで
$\mathscr P_{n,1}$の閉包に含まれる。
従って階数層の閉包が全ての既約成分を与え、その数は$\lfloor N/3\rfloor$である。
連結群$\mathrm{PGL}(V)$の商でも成分数は変わらない。
\end{proof}

\subsection{成分数の比較}
集合分割の指数母関数は
\begin{equation}\label{eq:pg-fermat-egf}
 \sum_{N\ge0}\frac{t^N}{N!}
 \left\{\begin{smallmatrix}N\\k\end{smallmatrix}\right\}_{\ge3}
 =\frac{(e^t-1-t-t^2/2)^k}{k!}.
\end{equation}
これと\eqref{eq:pg-fermat-stack-count}から次の値を得る。
\begin{center}
\small
\begin{tabular}{r|rrrr|rrr}
\toprule
$n$ & 階数$2$ & 階数$4$ & 階数$6$ & 階数$8$
& 固定空間 & 固定スタック & 相対スタック \\
\midrule
3 & 1 & -- & -- & -- & 1 & 1 & 1 \\
4 & 1 & 10 & -- & -- & 11 & 2 & 2 \\
5 & 1 & 35 & -- & -- & 36 & 2 & 2 \\
6 & 1 & 91 & -- & -- & 92 & 3 & 2 \\
7 & 1 & 210 & 280 & -- & 491 & 4 & 3 \\
8 & 1 & 456 & 2100 & -- & 2557 & 5 & 3 \\
9 & 1 & 957 & 10395 & -- & 11353 & 6 & 3 \\
10 & 1 & 1969 & 42735 & 15400 & 60105 & 9 & 4 \\
\bottomrule
\end{tabular}
\end{center}
中央の四列は固定Fermat空間の階数別成分数である。
例えば$N=9$では、階数$4$の$210$成分は$(3,6)$型の$84$個と
$(4,5)$型の$126$個からなり、固定スタックでは二つの成分になる。
一般の滑らかな3次超曲面の固定空間と固定スタックは、
系\ref{cor:pg-generic}によりともに既約である。

